% Abstract (150 words target)
% TODO: Complete after results are available

All algorithmic redistricting methods depend on a choice of atomic spatial units---typically census tracts, block groups, or blocks. This choice introduces the Modifiable Areal Unit Problem (MAUP): different spatial aggregations can yield different results even with identical data. We empirically test whether key findings from automated redistricting are sensitive to spatial resolution by comparing recursive bisection at three census geographic levels: tracts ($\sim$85,000 units), block groups ($\sim$240,000), and blocks ($\sim$11 million). [Results TBD]. This represents the first systematic MAUP analysis in algorithmic redistricting and establishes the generalizability---or lack thereof---of computational approaches to district design.
